%%%%%%%%%%%%%%%%%%%%%%%%%%%%%%%%%%%%%%%%%%%%%%%%%%%%%%%%%%%%%%%%%%%%%
%  sub/kr-2-framework.tex  —  2. 이론적 배경 (한글본)
%  영문본 sub/en-2-framework.tex 의 번역본.
%%%%%%%%%%%%%%%%%%%%%%%%%%%%%%%%%%%%%%%%%%%%%%%%%%%%%%%%%%%%%%%%%%%%%
\section{이론적 배경}\label{sec:framework-kr}

% 연구를 뒷받침하는 이론적 틀, 핵심 개념, 분석 렌즈를 제시합니다. 연구 문제와
% 방법(코딩 체계, 루브릭 등)을 이 틀에 연결하면 이후 결과 해석의 근거가 됩니다.

연구를 뒷받침하는 이론적 틀, 핵심 개념, 분석 렌즈를 제시하고, 연구 문제와 연구
방법을 이 틀에 연결합니다 \citep{examplebook2020}.
